\subsection{\texorpdfstring{SUMMABLE FAMILIES IN \(\mathbf{R}^n\)}{SUMMABLE FAMILIES IN R TO THE N}}

Since every point of \(\mathbf{R}^n\) has a countable fundamental system of neighbourhoods, a family \((\pmb{x}_i)\) of points of the additive group \(\mathbf{R}^n\) is summable only if the set of indices \(i\) such that \(\mathbf{x}_i \neq 0\) is countable (Chapter III, § 5, no. 2, Corollary to Proposition 1); hence, essentially, the study of summable families in \(R^{n}\) is reduced to that of summable sequences. Nevertheless, for the same reasons as were given in Chapter IV, § 7, in connection with summable families in R, we shall not impose any restriction, in what follows, on the cardinal of the index set.

PROPOSITION 1. A family \((\pmb{x}_t)_{t\in \mathbf{I}}\) of points \(\pmb{x}_t = (x_{t,k})_{1\leq k\leq n}\) of \(\mathbf{R}^n\) is summable if and only if each of the \(n\) families \((x_{t,k})_{t\in \mathbf{I}}\) of real numbers is summable in \(\mathbf{R}\) .

This follows from Chapter III, § 5, no. 4, Proposition 4.

The condition of Proposition 1 may be transformed as follows:

THEOREM 1. A family \((\pmb{x}_t)_{t\in \mathbf{I}}\) of points of \(\mathbf{R}^n\) is summable if and only if the family \((||\pmb{x}_t||)\) of Euclidean norms of the \(\pmb{x}_t\) is summable in \(\mathbf{R}\) .

This follows without trouble from Proposition 1, the condition for summability of a family of real numbers (Chapter IV, § 7, no. 2, Theorem 3), the inequalities

\[
\sup _ {1 \leqslant k \leqslant n} | x _ {t, k} | \leqslant | | x _ {t} | | \leqslant \sum_ {i = 1} ^ {n} | x _ {t, i} |,
\]

and the comparison principle (Chapter IV, § 7, no. 1, Theorem 2).

One can also proceed somewhat differently, by first establishing the following proposition:

PROPOSITION 2. If \((\pmb{x}_i)_{i\in \mathbf{I}}\) is any finite family of points in \(\mathbf{R}^n\) , then

\[
\sum_ {i \in \mathbf {I}} | | \boldsymbol {x} _ {i} | | \leqslant 2 n. \sup _ {\mathbf {J} \subset \mathbf {I}} \left| \left| \sum_ {i \in \mathbf {J}} \boldsymbol {x} _ {i} \right| \right|.\tag{1}
\]

For if \(\pmb{x}_i = (x_{ij})_{1\leqslant j\leqslant n}\) , we have \(||\pmb{x}_i|| \leqslant \sum_{j=1}^{n} |x_{ij}|\) , hence

\[
\sum_ {i \in \mathbf {I}} | | \boldsymbol {x} _ {i} | | \leqslant \sum_ {j = 1} ^ {n} \left(\sum_ {i \in \mathbf {I}} | x _ {i j} |\right).
\]

Now \(\sum_{i\in \mathbf{I}}|x_{ij}| = \sum_{i\in \mathbf{I}}x_{ij}^{+} + \sum_{i\in \mathbf{I}}x_{ij}^{-}\) , and since for every subset \(J\) of \(I\) we have

\[
- \sum_ {i \in \mathbf {I}} x _ {i j} ^ {-} \leqslant - \sum_ {i \in \mathbf {J}} x _ {i j} ^ {-} \leqslant \sum_ {i \in \mathbf {J}} x _ {i j} ^ {+} \leqslant \sum_ {i \in \mathbf {I}} x _ {i j} ^ {+},
\]

it follows that

\[
\sum_ {i \in \mathbf {I}} | x _ {i j} | \leqslant 2. \sup _ {\mathbf {J} \subset \mathbf {I}} \left| \sum_ {i \in \mathbf {J}} x _ {i j} \right|.
\]

But \(\left|\sum_{i\in J}x_{ij}\right| \leqslant \left\| \sum_{i\in J}\boldsymbol{x}_i\right\|\) , hence the inequality (1).

Now, Theorem 1 is equivalent to the following proposition (since \(\mathbf{R}^n\) is a complete group): the family \((\pmb{x}_t)\) satisfies Cauchy's criterion (Chapter III, § 5, no. 2, Theorem 1) if and only if the family \((||\pmb{x}_t||)\) also satisfies Cauchy's criterion. Now the triangle inequality shows that this condition is sufficient, and the inequality (1) shows that it is necessary.

Furthermore we have the inequality

\[
\left| \left| \sum_ {i} x _ {i} \right| \right| \leqslant \sum_ {i} | | x _ {i} | |,\tag{2}
\]

which comes by passing to the limit from the analogous inequality for finite partial sums.

COROLLARY. A family \((\pmb{x}_t)\) of points of \(\mathbf{R}^n\) is summable if and only if the set of finite partial sums of the family is bounded in \(\mathbf{R}^n\) .

By Theorem 1 and the triangle inequality, this condition is necessary; it is sufficient by the inequality (1) and Theorem 1.

PROPOSITION 3. Let \((\pmb{x}_{\lambda})_{\lambda \in \mathbb{L}}\) be a summable family of points of \(\mathbf{R}^m\) , \((y_{\mu})_{\mu \in \mathbb{M}}\) a summable family of points of \(\mathbf{R}^n\) , and let \(f\) be a bilinear mapping of \(\mathbf{R}^m \times \mathbf{R}^n\) into \(\mathbf{R}^p\) . Then the family \((f(x_{\lambda}, y_{\mu}))_{(\lambda, \mu) \in \mathbb{L} \times \mathbb{M}}\) is summable and we have

\[
\sum_ {(\lambda , \mu) \in \mathbf {L} \times \mathbf {M}} f \left(\boldsymbol {x} _ {\lambda}, \boldsymbol {y} _ {\mu}\right) = f \left(\sum_ {\lambda \in \mathbf {L}} \boldsymbol {x} _ {\lambda}, \sum_ {\mu \in \mathbf {M}} \boldsymbol {y} _ {\mu}\right).\tag{3}
\]

To show that the family \((f(x_{\lambda}, y_{\mu}))\) is summable, it is sufficient by Proposition 1 to establish that each of the p families formed by the coordinates of the points \(f(x_{\lambda}, y_{\mu})\) in \(R^{n}\) is summable: in other words, we can restrict ourselves to the case where f is a bilinear form; but for such a form f we have

\[
f (\boldsymbol {x}, \boldsymbol {y}) = \sum_ {i, j} a _ {i j} x _ {i} y _ {j},
\]

and therefore we are brought back to the case \(f(x, y) = x_i y_j\) , and in this case the result has already been proved (Chapter IV, § 7, no. 3, Proposition 1).

By specializing the function \(f\) we obtain in particular the following corollaries:

COROLLARY 1. If \((a_{\lambda})_{\lambda \in L}\) is a summable family of real numbers and if \((x_{\mu})_{\mu \in M}\) is a summable family of points of \(\mathbf{R}^n\) , then the family \((a_{\lambda}x_{\mu})_{(\lambda, \mu) \in L \times M}\) is summable and we have

\begin{enumerate}
\def\labelenumi{(\arabic{enumi})}
\setcounter{enumi}{3}
\tightlist
\item
\end{enumerate}

\[
\sum_ {(\lambda , \mu) \in \mathbf {L} \times \mathbf {M}} a _ {\lambda} x _ {\mu} = \left(\sum_ {\lambda \in \mathbf {L}} a _ {\lambda}\right) \left(\sum_ {\mu \in \mathbf {M}} x _ {\mu}\right).
\]

COROLLARY 2. If \((\pmb{x}_{\lambda})_{\lambda \in \mathbb{L}}\) and \((\pmb{y}_{\mu})_{\mu \in \mathbb{M}}\) are two summable families of points of \(\mathbf{R}^n\) , then the family \((\pmb{x}_{\lambda}|\pmb{y}_{\mu})\) (cf.~Chapter VI, §2, no. 2) is summable in \(\mathbf{R}\) , and we have

\[
\sum_ {(\lambda , \mu) \in \mathbf {L} \times \mathbf {M}} (x _ {\lambda} | y _ {\mu}) = \left(\sum_ {\lambda \in \mathbf {L}} x _ {\lambda} \mid \sum_ {\mu \in \mathbf {M}} y _ {\mu}\right).\tag{5}
\]

